\section{\centering Solutions}

\subsection{Solutions to T/F Questions}

\begin{table}[!h]
	\centering
	\begin{tabular}{c@{\hspace{2.5em}}cp{0.75\textwidth}}
		true & false & \\
		\specialrule{\lightrulewidth}{0.25em}{1.25em}
		\bubble & \dbubble & The quadratic form defined by \(Q(\vec{x}) = \vec{x}^TBB^T \vec{x}\) is always positive definite. \\[0.75em]
		\bubble & \dbubble &If \(\sigma_1, \sigma_2, \dots , \sigma_n\) are the singular values of a \(n \times n\) matrix \(A\) , then 
					    \(\det A = \sigma_1\sigma_2 \dots\sigma_n \). \\[0.75em]
		\dbubble & \bubble & Any matrix has a singular value decomposition.\\[0.75em]
		\bubble & \dbubble & If \(Q\) is a positive definite quadratic form with \(Q(\vec{x}) = \vec{x}^TA\vec{x}\) and \(A = A^T\), then the largest eigenvalue of \(A\) is the maximum value \(Q\) takes on in its domain.\\[0.5em]
		\dbubble & \bubble & Given \(A \in \mathbb{R}^{m \times n}\),  the largest singular value of \(A\) is the largest value \(Q(\vec{x}) = \|A\vec{x}\|\) takes on under the constraint \(\|\vec{x}\| = 1\).\\[0.5em]
		\dbubble & \bubble & For any nonzero matrix \(C \in \mathbb{R}^{m \times n}\), \(F(\vec{x}) = \|C\vec{x}\|^2\) defines a quadratic form.\\[0.5em]
		\dbubble & \bubble & If \(A \in \mathbb{R}^{m \times n}\) and the transformation \(\vec{x} \mapsto A\vec{x}\) is one-to-one, then 
					    it is impossibe for \(A\) to any singular values that equal to 0.\\[0.5em]
		\bottomrule
	\end{tabular}
\end{table}

\subsection{Solutions to Short Answer Questions}

\begin{enumerate}
	\item Make the change of variables \(x = Py\) using the orthogonal diagonalization of the symmetrix matrix corresponding to \(Q(\vec{x})\), and you get
		\(Q'(\vec{y}) = 2y_1^2 + 4y_2^2\).
		\begin{enumerate}
			\item 4.
			\item \(\vec{x} = \pm\begin{bmatrix}0 \\ 1 \\ 0\end{bmatrix}\).
			\item This is the second largest eigenvalue of the matrix we found above, which is 2.
		\end{enumerate}
	 
	\item SVD: \[\begin{bmatrix}1 & -1 \\ 0 & 0 \\ 1 & -1\end{bmatrix} = \underbrace{\begin{bmatrix}1/\sqrt{2} & 0  & -1/\sqrt{2} \\ 0 & 1 & 0 \\ 1/\sqrt{2} & 0 & 1/\sqrt{2}\end{bmatrix}}_U \underbrace{\begin{bmatrix}2 & 0 \\ 0 & 0 \\ 0 & 0\end{bmatrix}}_\Sigma\underbrace{\begin{bmatrix}1/\sqrt{2} & -1/\sqrt{2} \\ 1/\sqrt{2} & 1/\sqrt{2}\end{bmatrix}}_{V^T}.\]
		\begin{enumerate}
			\item \(\left\{ \begin{bmatrix}1/\sqrt{2} \\ 0 \\ 1/\sqrt{2}\end{bmatrix} \right\}\).\footnote{In general, if \(\rank A = r\), then the first \(r\)
				 columns of \(U\) are form an orthonormal basis for the column space. The first \(r\) columns of \(V\) would also form an orthonormal
				 basis for \(\rowspace A\).}
			\item \(\left\{ \begin{bmatrix}1/\sqrt{2} \\ -1/\sqrt{2}\end{bmatrix} \right\}\).
			\item \(\left\{ \begin{bmatrix}1/\sqrt{2} \\ 1/\sqrt{2}\end{bmatrix} \right\}\).
			\item \(\left\{ \begin{bmatrix}0 \\ 1 \\ 0\end{bmatrix},  \begin{bmatrix}-1/\sqrt{2} \\ 0 \\ 1/\sqrt{2}\end{bmatrix}\right\}\).
			\item \(\rank A = 1\).
		\end{enumerate}
	\item \(\max_{\|\vec{x}\| = 1} \|A\vec{x}\| = \sigma_1 = 2\). This value is attained when 
		 \(\vec{x} = \pm\begin{bmatrix}1/\sqrt{2} \\ -1/\sqrt{2}\end{bmatrix}\).
	\item \(\min_{\|\vec{x}\| = 1} \|A\vec{x}\| = \sigma_2 = 0\). This value is attained when 
		\(\vec{x} = \pm\begin{bmatrix}1/\sqrt{2} \\ 1/\sqrt{2}\end{bmatrix}\).

\end{enumerate} 
